\appendix
\section*{Appendix}
\addcontentsline{toc}{section}{Appendix}
\renewcommand{\thesubsection}{\Alph{subsection}}
\setcounter{subsection}{0}

This appendix collects the material that supports the main text without interrupting its argument. It
states the prediction algorithm (\ref{app:algorithm}) and the study's positionality (\ref{app:positionality}),
gives the mathematical derivations behind the factor and prediction models (\ref{app:math}), the full
coverage audit behind the inclusion rule (\ref{app:coverage}), the tuning grid (\ref{app:hyper}), the
robustness battery (\ref{app:robust}), the reproducibility and provenance record (\ref{app:provenance}),
and the supplementary figures referenced from the results (\ref{app:figures}). Each subsection opens with
a one-line statement of what it supports and where the main text relies on it.

\subsection{Leave-one-benchmark-out procedure}\label{app:algorithm}
\emph{Supports Section~\ref{sec:task2}, which describes the protocol in words.}
\begin{figure}[H]
\centering
\fbox{\parbox{0.95\linewidth}{
\small
\textbf{Procedure}\quad Leave-one-benchmark-out predictive validity with nested cross-validation\\[2pt]
\hrule
\vspace{3pt}
\textbf{Input:} standardised score matrix $Z$ (models $\times$ benchmarks), economic targets $T$, learners $L$, rungs $R$.\\[2pt]
\textbf{for} each target $t \in T$ \textbf{do}\\
\hspace*{1em} partition models into 5 outer folds\\
\hspace*{1em} \textbf{for} each outer fold $(\mathrm{train}, \mathrm{test})$ \textbf{do}\\
\hspace*{2em} fit standardisation and EFA on $\mathrm{train}$ only; form Thurstone weights $W = R^{-1}\Lambda$\\
\hspace*{2em} project $\mathrm{test}$ models through $W$; orient factors by train mean index\\
\hspace*{2em} \textbf{for} each rung $r \in R$ and learner $\ell \in L$ \textbf{do}\\
\hspace*{3em} tune $\ell$ by inner 5-fold grid search on $\mathrm{train}$; predict $\mathrm{test}$, store out-of-fold error\\
\hspace*{1em} select best learner per rung by pooled out-of-fold RMSE\\
compute $\Delta\mathrm{MSE} = \mathrm{MSE}(\text{rung ii}) - \mathrm{MSE}(\text{rung iv})$; bootstrap $B=2000$ for 95\% CI\\[2pt]
\textbf{Output:} per-target and pooled $\Delta\mathrm{MSE}$ with confidence intervals (H4).
}}
\caption[Leave-one-benchmark-out prediction procedure]{The leave-one-benchmark-out procedure with leakage-safe factor projection and nested tuning.}
\label{alg:lobo}
\end{figure}

\subsection{Positionality, directionality, and threats to validity}\label{app:positionality}
\emph{Supports Section~\ref{sec:hypotheses}, which summarises these commitments in one sentence.}
Because a factor solution is identified only up to sign, the direction of every factor is a choice we
state explicitly. We orient each factor to correlate positively with the model's mean standardised
score, so that ``higher'' always means ``more capable'' and every loading in Table~\ref{tab:loadings}
reads in the same direction. The point is not cosmetic, since an un-oriented solution could report the
same economic factor with flipped signs and invite the opposite interpretation. The direction of the
residualisation is equally deliberate. We regress scores on release date and analyse the residuals,
which asks how much benchmark co-movement remains once the release date is known, and we do not run the
reverse regression, because we make no claim that the general factor causes the calendar. Finally we fix
the sign of the predictive contrast in advance, defining $\Delta\mathrm{MSE}$ as the single-index error
minus the $k$-factor error so that a positive value favours the richer representation, and pre-register
that H4 requires this positive interval to exclude zero rather than testing the reverse direction.

Two threats to validity follow from the vantage point of the study. First, we observe only the models
that vendors chose to submit to a public leaderboard, so the sample is a convenience sample of frontier
systems and not a random draw, and the claims generalise to models on this leaderboard rather than to
all models. Second, the economic benchmarks are newer and sparser, so any economic-specific structure is
estimated from fewer observations and is more sensitive to the particular set of models that happen to
carry those scores, a limitation we quantify with the wide bootstrap intervals in Table~\ref{tab:h4}
rather than hide behind point estimates.

\subsection{Mathematical derivations}\label{app:math}
\emph{Supports Sections~\ref{sec:task1} and \ref{sec:task2}, which state the factor and prediction models
in words and defer their formal statement to here.}

\medskip
\noindent\textbf{The exploratory factor-analysis model.}
Let $\mathbf{z}\in\mathbb{R}^{p}$ be the vector of $p$ per-benchmark z-standardised scores for a model.
The common-factor model represents it as a linear combination of $k$ latent factors
$\mathbf{f}\in\mathbb{R}^{k}$ plus a benchmark-specific uniqueness $\boldsymbol{\varepsilon}$,

\[
  \mathbf{z} = \Lambda \mathbf{f} + \boldsymbol{\varepsilon},
\]

where $\Lambda\in\mathbb{R}^{p\times k}$ is the loading matrix. Assuming
$\mathrm{Cov}(\mathbf{f})=\Phi$, $\mathrm{Cov}(\boldsymbol{\varepsilon})=\Psi$ diagonal, and
$\mathrm{Cov}(\mathbf{f},\boldsymbol{\varepsilon})=\mathbf{0}$, the model implies the covariance
structure

\[
  \Sigma = \Lambda \Phi \Lambda^{\top} + \Psi .
\]

\medskip
\noindent
The parameters are estimated by maximum likelihood under multivariate normality. The dominance statistic
for H1 is computed on the unrotated solution, where the factors are orthogonal ($\Phi=I$) and each
factor's contribution to common variance is its sum of squared loadings, $\sum_{j=1}^{p}\lambda_{jm}^2$,
the $m$-th diagonal entry of $\Lambda^{\top}\Lambda$. The first factor's share of common variance is
therefore
\[
  s_1 \;=\; \frac{\sum_{j=1}^{p}\lambda_{j1}^2}{\sum_{m=1}^{k}\sum_{j=1}^{p}\lambda_{jm}^2}
        \;=\; \frac{(\Lambda^{\top}\Lambda)_{11}}{\operatorname{tr}(\Lambda^{\top}\Lambda)} .
\]
On the complete-case grid the three factors carry sums of squared loadings $7.79$, $2.34$, and $0.32$, so
$s_1 = 7.79/10.45 = 74.5\%$. This is a factor-level share, distinct from a single benchmark's communality
$\sum_{m}\lambda_{jm}^2$, which is a diagonal entry of $\Lambda\Lambda^{\top}$.

\medskip
\noindent\textbf{Oblique rotation and factor scores.}
The maximum-likelihood solution is identified only up to rotation, so we apply an oblique oblimin
rotation, which minimises a quartimin-type criterion on the loadings while allowing $\Phi\neq I$. Factor
scores are then recovered by the Thurstone regression method,

\[
  \mathbf{W} = R^{-1}\Lambda, \qquad \hat{\mathbf{f}} = \mathbf{W}^{\top}\mathbf{z},
\]

where $R$ is the sample correlation matrix. In the prediction task $\mathbf{W}$ is estimated on the
training fold only and applied to the held-out fold, which is the leakage-safe projection referenced in
Section~\ref{sec:task2}.\label{app:scores}

\medskip
\noindent\textbf{Date and scale residualisation.}
To separate general capability from a shared time trend, each standardised benchmark $z_j$ is regressed
on model release date $d$ (and, in a subsample, on log training compute),

\[
  z_j = \beta_{0j} + \beta_{1j}\, d + r_j ,
\]

and the factor model is re-fitted on the residuals $r_j$. The change in the first factor's variance
share between the raw and residualised fits, $74.5\%\rightarrow 59.6\%$, is the H2(ii) statistic.

\medskip
\noindent\textbf{Nested leave-one-benchmark-out estimator.}
For an economic target $t$, the out-of-fold error of rung $r$ is

\[
  \mathrm{MSE}_r = \frac{1}{n}\sum_{i=1}^{n}\bigl(z_{it} - \hat{z}_{it}^{(-\kappa(i))}\bigr)^{2},
\]

where $\hat{z}_{it}^{(-\kappa(i))}$ is the prediction for model $i$ from a learner trained on the outer
folds excluding the fold $\kappa(i)$ containing $i$, with hyperparameters chosen by an inner five-fold
grid search on those training folds. The predictive-validity statistic contrasts the single-index rung
against the $k$-factor rung,

\[
  \Delta\mathrm{MSE} = \mathrm{MSE}_{\text{ii}} - \mathrm{MSE}_{\text{iv}} .
\]

\medskip
\noindent\textbf{Bootstrap confidence interval.}
A 95\% interval for $\Delta\mathrm{MSE}$ is obtained by resampling the paired out-of-fold squared-error
vectors with replacement $B=2000$ times and taking the 2.5th and 97.5th percentiles of the resampled
differences. H4 is supported when this interval excludes zero.

\subsection{Benchmark battery}\label{app:battery}
\emph{Supports Section~\ref{sec:battery}, which names the battery and defers the full description here.}
\begin{table}[tb]
\centering
\caption[The twelve-benchmark battery]{\textbf{The twelve-benchmark battery, grouped by capability block,} with the construct each
targets and its coverage in the analysis snapshot (models scored, of 421 retained). Higher is better for
all; GDPval is an Elo rating, the others accuracy-type scores.}
\label{tab:battery}
\small
\begin{tabular}{@{}llr@{}}
\toprule
Benchmark & Construct measured & Coverage \\
\midrule
\multicolumn{3}{@{}l}{\textit{Economic}}\\
GDPval (Elo) & Paid professional task quality & 112 \\
Terminal-Bench v2.1 & Terminal / agentic task completion & 121 \\
$\tau^3$-Banking & Multi-step banking workflows & 112 \\
$\tau^2$-Bench & Tool-use agentic tasks & 413 \\
\multicolumn{3}{@{}l}{\textit{Academic}}\\
GPQA Diamond & Graduate-level science QA & 421 \\
HLE & Humanity's Last Exam (frontier knowledge) & 421 \\
AA-Omniscience & Broad knowledge, hallucination-penalised & 418 \\
\multicolumn{3}{@{}l}{\textit{Scientific coding}}\\
SciCode & Research-level scientific coding & 421 \\
CritPt & Physics problem solving & 419 \\
Terminal-Bench Hard & Hard terminal coding tasks & 413 \\
\multicolumn{3}{@{}l}{\textit{Long-context / instruction}}\\
AA-LCR & Long-context reasoning & 420 \\
IFBench & Instruction following & 417 \\
\bottomrule
\end{tabular}
\end{table}

\subsection{Full coverage audit}\label{app:coverage}
\emph{Supports the inclusion rule of Section~\ref{sec:inclusion};} Table~\ref{tab:coverage} reports every
candidate benchmark on the raw snapshot and the decision applied to it.

\begin{table}[htbp]
\centering
\caption[Full candidate-benchmark coverage audit]{Full candidate-benchmark coverage on the raw 548-model snapshot, with the inclusion decision.
Twelve benchmarks are primary, MMMU-Pro is demoted to sensitivity-only for sparse and multimodal
coverage, and APEX-Agents is dropped below the sixty-model floor.}
\label{tab:coverage}
\small
\begin{tabular}{@{}llrl@{}}
\toprule
Benchmark & Block & Models scored & Role \\
\midrule
GPQA Diamond & Academic & 513 & primary \\
HLE & Academic & 509 & primary \\
AA-Omniscience & Academic & 418 & primary \\
MMMU-Pro & Academic & 204 & sensitivity only \\
$\tau^2$-Bench & Economic & 428 & primary \\
Terminal-Bench v2.1 & Economic & 121 & primary \\
GDPval (Elo) & Economic & 117 & primary \\
$\tau^3$-Banking & Economic & 112 & primary \\
APEX-Agents & Economic & 26 & dropped (sparse) \\
IFBench & Instruction-following & 437 & primary \\
AA-LCR & Long-context & 441 & primary \\
SciCode & Scientific-coding & 507 & primary \\
CritPt & Scientific-coding & 422 & primary \\
Terminal-Bench Hard & Scientific-coding & 421 & primary \\
\bottomrule
\end{tabular}
\end{table}

\subsection{Hyperparameter grid and selections}\label{app:hyper}
\emph{Supports Section~\ref{sec:task2} and the robustness claim of Section~\ref{sec:robustness}} that the
H4 verdict does not hinge on a single tuning choice. The inner grid searched ridge and elastic-net
$\alpha\in\{0.03,0.1,0.3,1,3\}$ with elastic-net $\ell_1$ ratio in $\{0.2,0.5,0.8\}$, random-forest
depth in $\{4,8,\text{none}\}$ with $\{200,300\}$ trees, and gradient-boosting depth in $\{2,3\}$ with
learning rate in $\{0.05,0.1\}$. Table~\ref{tab:hyper} reports the modal selection per rung.

\begin{table}[htbp]
\centering
\caption[Modal selected hyperparameters]{Modal hyperparameters selected by the inner five-fold grid search, reported as the plurality
choice across all twelve leave-one-benchmark-out targets with the agreement count. The H4 verdict rests
on fold-averaged out-of-fold errors, not on any single setting.}
\label{tab:hyper}
\small
\begin{tabular}{@{}lllc@{}}
\toprule
Rung & Learner & Modal hyperparameters & Agreement \\
\midrule
(i) timing / scale & random forest & \texttt{max\_depth}=4, \texttt{min\_samples\_leaf}=3, \texttt{n\_estimators}=300 & 8/12 \\
(ii) mean index & ridge & \texttt{alpha}=0.03 & 6/12 \\
(iii) first factor & elastic net & \texttt{alpha}=0.03, \texttt{l1\_ratio}=0.2 & 5/12 \\
(iv) $k$ factors & ridge & \texttt{alpha}=0.03 & 6/12 \\
(v) $k$ factors + covariates & ridge & \texttt{alpha}=1 & 6/12 \\
\bottomrule
\end{tabular}
\end{table}

\subsection{Deviations from the pre-registered design}\label{app:deviations}
\emph{Supports the reproducibility claims throughout.} A study that invokes pre-registration owes the
reader an explicit account of where the delivery departed from the plan. Four departures matter. First,
the registered primary unit was one row per base model at its best default configuration; the analysis
instead uses all 421 configurations, with the deduplicated base-model dataset reported below as the
registered R1 check rather than as the primary. Second, the registered factor-count rule is parallel
analysis, which retains one factor; the three-factor solution that exhibits economic separation is a
pre-declared exploratory over-extraction, and H3 is reported as not supported under the registration in
consequence. Third, the taxonomy places Terminal-Bench Hard in scientific coding whereas the
registration grouped it with Terminal-Bench v2.1 as economic; the residualised loadings show Hard at
0.82 on the economic factor, so the assignment is genuinely contestable and the split is disclosed
rather than silently adopted. Fourth, the delivered robustness checks are not the registered suite, and
we label them distinctly below instead of renumbering a different battery onto the registered names.

\subsection{Robustness checks}\label{app:robust}
\emph{Supports Section~\ref{sec:robustness}.} We first report the one registered check that bears
directly on the deduplication concern, then the additional checks actually delivered. \textbf{Registered
R1 (configuration versus deduplicated).} We collapse the 421 configurations to one row per base model,
keeping the best default configuration as the registration specified, operationalised as the
highest-intelligence-index row for each base model after stripping reasoning, effort, and preview
suffixes from the model name. This leaves 320 base models, of which 89 carry the complete battery. On this deduplicated grid the
first-factor share rises to 90.2\% and the date-residualisation drop rises to 24.1 points, and the four
economic loadings remain above 0.40, so deduplication strengthens H1 and H2(ii) and leaves the
exploratory economic pattern intact. The configuration-level analysis is therefore the conservative
choice for the structural headline. We also re-run the H4 predictive test on this deduplicated grid,
where the bootstrap resampling unit is a distinct base model rather than a configuration. The pooled
economic-block improvement of the $k$-factor representation over the single index is $+0.038$ with a 95\%
interval of $[+0.020, +0.056]$, matching the configuration-level $+0.037$ and again excluding zero, so
the predictive gain does not depend on treating near-duplicate configurations as exchangeable. \textbf{Delivered checks C1--C6.} (C1) A rank-based factor analysis
raises the first-factor share to 94.5\%, and (C2) a logit transform of the bounded accuracy benchmarks
raises it to 91.8\%, so H1 is not a Pearson artefact. (C3) On the 58-model compute-known subsample, the
date-only residualisation drops the first-factor share by 16.5 points and adding log-compute jointly
gives a smaller 9.3-point drop, so compute does not deepen the temporal correction beyond date on the
same models. (C4) The three clustering algorithms agree on $k=2$. (C5) The predictor-ladder ordering
holds on both the economic block and all twelve targets. (C6) The dominant-factor date fit is robust to
the link function: a logistic fit gives $R^2=0.505$ against an ordinary-least-squares $R^2=0.477$, both
clearing 0.30 with the first factor ranked strongest, so H2(i) does not depend on that modelling choice.
(C7) The $k$-factor predictive advantage does not depend on the exploratory factor count. Sweeping the
rung-(iv) predictor over $k \in \{2,3,4,5\}$, the improvement over the single index is positive at every
count with a bootstrap interval excluding zero: $+0.021$ at $k=2$, $+0.037$ at the $k=3$ we use, $+0.058$
at the criterion's $k=4$, and $+0.059$ at $k=5$, with pooled test $R^2$ rising from 0.79 to 0.83 across
the range. The gain increases with $k$, so using the $k=3$ we adopt is conservative for this test
relative to the criterion's $k=4$.
The registered checks not run are R3 (FIML and MICE missingness sensitivity), R4 (replication of the
single-factor foil then correction), R5 (reasoning-effort variants as a second scale axis), and R6 (the
taxonomy swap test); each was deferred for space and is the natural next step given the deviations above,
particularly R5 and R6, which bear on the effort and taxonomy departures.

\subsection{Data provenance and reproducibility}\label{app:provenance}
\emph{Supports Section~\ref{sec:source} and the reproducibility claims throughout.} The primary data set
is the Artificial Analysis model leaderboard \cite{aa2026}, captured 2026-07-06 as a raw HTML snapshot
and pinned by the SHA-256 hash recorded in the shipped manifest
(\texttt{aa\_snapshot\_manifest.json}),
\texttt{\seqsplit{6f19f8f08befaaf14cb2d9952e379404c39f246b52d7ae4d1c2aa44795404de6}}. The Data API endpoint returned
HTTP 401 to our unauthenticated request, so the public leaderboard page was captured and the exact bytes
stored, which makes every downstream number reproducible regardless of subsequent live-page drift. The
secondary data set is Epoch AI Notable AI Models \cite{epoch2026}, released under CC-BY 4.0, used only
for the scale robustness check. Both snapshots and every result table travel with the submission bundle.
The complete pipeline is in a single notebook, \texttt{analysis.ipynb}, which regenerates all
figures and tables from the pinned inputs, and the full code and data will be released at
\url{https://github.com/louisyzhu/frontier-ai-economic-validity}.

\subsection{Supplementary figures}\label{app:figures}
\emph{Supports Sections~\ref{sec:inclusion}, \ref{sec:h1}, and \ref{sec:clustering}.}
Figure~\ref{fig:density} documents the coverage audit behind the inclusion rule, Figure~\ref{fig:clusters}
shows the model clustering, Figure~\ref{fig:diag} the GDPval regression diagnostics, Figure~\ref{fig:scree} gives the dimensionality evidence behind H1, and
Figure~\ref{fig:dendro} the benchmark-level clustering behind Section~\ref{sec:clustering}.

\begin{figure}[htbp]
  \centering
  \includegraphics[width=\linewidth]{fig3_clusters}
  \caption[Model clustering into two capability tiers]{\textbf{Model clustering on the dense grid ($n=409$).} (a) Models in the space of the first
  two factors, coloured by the two-cluster agglomerative solution; the factor axes are the dense-grid
  solution fitted largely without the economic columns, so they index general capability tiers rather
  than economic type. (b) Silhouette scores by $k$ for all three algorithms agree on $k=2$.}
  \label{fig:clusters}
\end{figure}

\begin{figure}[htbp]
  \centering
  \includegraphics[width=\linewidth]{appxfig_density}
  \caption[Coverage and pairwise-overlap audit]{\textbf{Coverage audit behind the inclusion rule.}
  (a) Number of models carrying a score on each candidate benchmark, coloured by capability block, with
  the sixty-model inclusion floor marked. Four of the five economic benchmarks are retained; three of
  them (Terminal-Bench v2.1, GDPval, and $\tau^3$-Banking) are scored on only 112 to 121 models, whereas
  $\tau^2$-Bench and the academic core are each scored on four hundred or more. APEX-Agents, with 26
  models, falls below the floor and is dropped. (b) Pairwise counts of models scored
  on both benchmarks in a pair, which set the effective sample size for every correlation in the study.
  This is the evidence for the two-grid design of Section~\ref{sec:inclusion}.}
  \label{fig:density}
\end{figure}

\begin{figure}[htbp]
  \centering
  \includegraphics[width=\linewidth]{suppl_eda_distributions}
  \caption[Univariate score distributions and skewness]{\textbf{Univariate benchmark score
  distributions.} (a) Raw score histograms for the twelve benchmarks, coloured by capability block;
  $\gamma$ is the Fisher skewness. (b) Skewness ordered across benchmarks. Most benchmarks are bounded
  accuracies in $[0,1]$, and several are strongly right-skewed (CritPt $\gamma=3.8$, HLE $\gamma=1.5$,
  $\tau^3$-Banking $\gamma=0.9$), with scores concentrated near the floor. This boundedness motivates the
  logit-transform robustness check of Section~\ref{sec:robustness}. GDPval (Elo) and AA-Omniscience are
  on unbounded scales and are closer to symmetric.}
  \label{fig:eda_dist}
\end{figure}

\begin{figure}[htbp]
  \centering
  \includegraphics[width=\linewidth]{fig5_diagnostics}
  \caption[Regression diagnostics for the GDPval model]{\textbf{Regression diagnostics for the GDPval $k$-factor ridge model, out-of-fold.}
  (a) Predicted versus observed, on the identity line with $R^2=0.88$. (b) Residuals versus predicted
  show no systematic trend. (c) The eight largest residuals, named; these are interpretable outliers
  rather than evidence of misspecification.}
  \label{fig:diag}
\end{figure}

\begin{figure}[htbp]
  \centering
  \includegraphics[width=\linewidth]{suppl_scree_parallel}
  \caption[Scree plot and Horn parallel analysis]{\textbf{Dimensionality of the battery.} (a) The scree plot of principal-component variance
  shares drops sharply after the first component (79.4\%). (b) Horn parallel analysis retains a single
  factor, since only the first observed eigenvalue (9.53) exceeds its random-data 95th percentile (1.79),
  while the second (0.91) already falls below (1.57). Values are computed in the notebook and stored in
  \texttt{parallel\_analysis.csv}.}
  \label{fig:scree}
\end{figure}

\begin{figure}[htbp]
  \centering
  \includegraphics[width=\linewidth]{suppl_benchmark_dendrogram}
  \caption[Benchmark clustering dendrogram]{\textbf{Benchmark clustering by correlation distance} ($1-\rho$, average linkage). Ten of the
  twelve benchmarks, including three of four economic benchmarks, fall in one block, with $\tau^2$-Bench
  and IFBench branching separately.}
  \label{fig:dendro}
\end{figure}